\section{Problem Formulation}

Let $S_0$ be the pre-deletion system and $S_-$ be the post-deletion system. A system may include a base model, retriever, vector index, reranker, prompt template, adapter, cache, and operational logs. A deletion request targets a set of data units $D_{\mathrm{del}}=\{z_1,\ldots,z_k\}$. Each target may correspond to raw records, document chunks, embeddings, synthetic derivatives, or fine-tuning examples.

The counterfactual gold standard is a clean reference system $S_{\mathrm{clean}}(D_{\mathrm{del}})$ built as if the deletion targets had never entered any pipeline component: no raw record, chunk, embedding, cache item, synthetic derivative, adapter update, evaluation trace, or summary is produced from $D_{\mathrm{del}}$. Exact construction of $S_{\mathrm{clean}}$ may require retraining, rebuilding indexes, regenerating synthetic data, and purging caches. When that construction is infeasible, it remains the formal oracle that defines deletion success; practical audits estimate distance to this oracle.

We model data dependence with a provenance graph $G=(V,E)$. Nodes include deletion targets and downstream artifacts. An edge $(u,v)$ means artifact $v$ was produced from or influenced by $u$. Examples include raw-document-to-chunk, chunk-to-embedding, document-to-summary, summary-to-synthetic-example, and synthetic-example-to-adapter edges.

\textbf{Behavioral-deletion guarantee.}
Let $\mathcal{Q}(D_{\mathrm{del}},G)$ be a probe class and let $O_S(q)$ denote the observable behavior of system $S$ on probe $q$, including generated answer, retrieved context identifiers, citations, and available metadata. For a bounded discrepancy $d(\cdot,\cdot)\in[0,1]$, define
\[
\Delta_{\mathcal{Q}}(S_-,S_{\mathrm{clean}})
=
\sup_{q\in\mathcal{Q}}
d\!\left(P(O_{S_-}(q)),P(O_{S_{\mathrm{clean}}}(q))\right).
\]
We say $S_-$ satisfies $(\varepsilon,\delta)$ behavioral deletion for $D_{\mathrm{del}}$ over $\mathcal{Q}$ if, with probability at least $1-\delta$ over audit sampling and system randomness,
\[
\Delta_{\mathcal{Q}}(S_-,S_{\mathrm{clean}})\leq \varepsilon.
\]
This definition lifts counterfactual indistinguishability from model-centric unlearning to compositional AI pipelines: a post-deletion system must behave like the system that would have existed had the record never affected retrieval, caches, synthetic derivatives, or adapters.

The audit problem is to test whether $S_-$ is behaviorally indistinguishable from $S_{\mathrm{clean}}$ while preserving utility on a retain distribution $Q_{\mathrm{ret}}$. We do not require the auditor to inspect model weights. Instead, the auditor can submit black-box probes and observe generated answers, retrieved contexts, citations, and available metadata. The six evidence channels used later are estimators and diagnostics for this counterfactual distance; they are not the definition of deletion failure.

An audit returns one of three decisions. \emph{Pass} means the floor-normalized excess counterfactual risk is below the declared operating tolerance and retain utility is acceptable. \emph{Fail} means the post-deletion system is distinguishable from the clean counterfactual or a diagnostic channel localizes a violation. \emph{Escalate} means the evidence is inconclusive, the clean reference is unavailable, or retain utility has collapsed.

\textbf{Threat model.}
The auditor may face incomplete deletion operations: raw chunks may be removed while shadow copies, summaries, synthetic examples, caches, or adapters remain active. The baseline attacker class is black-box or gray-box: the auditor queries the post-deletion application and, in gray-box RAG settings, inspects retrieved context identifiers and citations, but does not assume access to model weights or training gradients. A white-box self-influence adversary with access to gradients, adapter weights, or internal memorization scores is outside the guarantee boundary by assumption; such tools can be added as privileged evidence. We do not claim that black-box auditing proves absolute absence of influence under all prompts; the objective is bounded evidence of counterfactual behavioral indistinguishability over a declared, adversary-aware probe class.
